% MT067 / explanation — canonical editable content.
导数\(dy/dx\)表示\(x\)在当前位置增加一点时，\(y\)变化得多快；例如\(y=3x\)的导数为3。多个变量参与计算时，\(\partial L/\partial a\)表示暂时固定其他变量，只看\(a\)改变对\(L\)的影响，称为偏导数。


链式法则适用于“一个计算结果又被送入下一步”：最终结果对输入的变化率，等于沿路径各局部变化率的乘积；若输入通过多条路径影响结果，则把各条贡献相加。先按(a)练一条标量计算链，再用于(b)的网络。


\textbf{(a) 把长表达式拆成短操作。} 这里用到三条基本求导规则：\((u^p)^{\prime}=pu^{p-1}\)、\((e^u)^{\prime}=e^u\)、\((\tan u)^{\prime}=1/\cos^2u\)，撇号都表示对各自输入\(u\)求导；三角函数角度按弧度计算。复合函数还要再乘内层输入的导数。定义
\[
s=3x^2,\quad r=s^{1/3},\quad t=5x,\quad v=\tan t,
\quad u=r+v,\quad q=e^u,\quad y=(1+q)^{-1}.
\]


从外向内求导。倒数层给\(\partial y/\partial q=-(1+q)^{-2}\)，指数层给\(\partial q/\partial u=e^u\)，所以
\[
\frac{dy}{du}=-\frac{e^u}{(1+e^u)^2}.
\]


\(u\)有两条分支，两条导数要相加。立方根分支先对\(s\)求导，再乘\(s\)对\(x\)的导数：



\[
\frac{dr}{dx}=\frac13s^{-2/3}\cdot6x
=\frac{2x}{(3x^2)^{2/3}}.
\]


另一分支\(\tan t\)的导数为\(\sec^2t=1/\cos^2t\)，再乘\(dt/dx=5\)，得\(dv/dx=5\sec^2(5x)\)。于是
\[
\frac{du}{dx}=\frac{2x}{(3x^2)^{2/3}}+5\sec^2(5x).
\]


把它乘上\(dy/du\)即为英文答案中的最终式。最容易漏的是内层的6x、5，以及外层倒数带来的负号。\(x=0\)时立方根复合项有尖点，普通导数不适用；\(\cos(5x)=0\)时tan无定义，也必须排除。


\textbf{(b) 先讲前向，每个变量都有尺寸。} 输入\(x\in\mathbb R^m\)；\(A\)有\(m\)列，每列\(a_i\)是第\(i\)个隐藏单元的输入权重。选双曲正切tanh作为隐藏激活：\(\tanh(t)=(e^t-e^{-t})/(e^t+e^{-t})\)，它把加权和压到\((-1,1)\)；选线性输出，即最后一层直接输出加权和。这两项都是原题允许自定的部分。前向为
\[
g_{1i}=a_i^Tx=\sum_{k=1}^m A_{ki}x_k,\qquad z_i=\tanh(g_{1i}).
\]


因此\(g_1,z\)各有\(m\)个数。\(W\)有\(n\)列，每列\(w_j\in\mathbb R^m\)连接所有隐藏单元到第\(j\)个输出：
\[
g_{2j}=w_j^Tz=\sum_{i=1}^m W_{ij}z_i,\qquad f_j=g_{2j}.
\]


所以\(g_2,f\)各有\(n\)个数。损失\(L(f,y)\)比较输出与目标；题目没有指定损失，不能擅自把所有后续结论限定为平方损失。


\textbf{从输出往回传。} 记\(\delta_{2j}=\partial L/\partial f_j\)。由于输出是恒等映射，也有\(\partial L/\partial g_{2j}=\delta_{2j}\)。一个输出权重\(W_{ij}\)只在\(f_j\)中乘\(z_i\)，所以
\[
\frac{\partial L}{\partial W_{ij}}=\delta_{2j}z_i,
\qquad \frac{\partial L}{\partial w_j}=z\,\delta_{2j}.
\]


第\(i\)个隐藏激活\(z_i\)会影响所有\(n\)个输出，因此必须把分支相加：
\[
\frac{\partial L}{\partial z_i}=\sum_{j=1}^{n}W_{ij}\delta_{2j}.
\]


tanh的局部导数为\(1-\tanh^2(g_{1i})=1-z_i^2\)，再乘它得到隐藏输入的梯度：
\[
\delta_{1i}:=\frac{\partial L}{\partial g_{1i}}
=\left(\sum_{j=1}^{n}W_{ij}\delta_{2j}\right)(1-z_i^2).
\]


最后\(g_{1i}=a_i^Tx\)，对\(a_i\)的各分量求导分别得到对应输入\(x_k\)，所以
\[
\frac{\partial L}{\partial a_i}=x\,\delta_{1i}.
\]


把逐分量结果收成向量，就是\(\delta_1=(W\delta_2)\odot(1-z^2)\)。\(W\)为\(m\times n\)，乘\(n\)维\(\delta_2\)得到\(m\)维，再与\(m\)维激活导数逐元素相乘；\(\odot\)不是矩阵乘法。按题目所列顺序，三个量依次为\(x\delta_{1i}\)、\(\delta_1\)和\(z\delta_{2j}\)，每个都是\(m\)维。


\textbf{用一个小数值检查方向。} 额外选择\(L=\frac12(f-y)^2\)，令\(m=n=1,x=1,A=0,W=2,y=1\)。前向\(g_1=z=f=0\)，因此\(\delta_2=f-y=-1\)，\(\delta_1=2(-1)(1-0)=-2\)。所以\(\partial L/\partial A=-2\)，\(\partial L/\partial W=z\delta_2=0\)。当前隐藏输出为0，微调输出权重暂时无效；增大\(A\)却可先产生正的隐藏输出，使\(f\)靠近目标1，这与负梯度更新方向一致。


选读：完整矩阵梯度为\(\nabla_A L=x\delta_1^T\)和\(\nabla_W L=z\delta_2^T\)，分别为\(m\times m\)、\(m\times n\)，与参数尺寸完全一致。若改成其他输出激活或损失，需要重新写对应局部导数，再沿同样路径相乘、相加。


\textbf{出处：} 2025A期中 Q13，原卷第10页（10分）


\textbf{答案依据：} 2025A期中 Q13 印刷题干；本题解为我们的推导，不采用扫描件的学生作答。
